\begin{abstract}
Advancements like BitNet~\cite{bitnet} signal a transformative shift towards 1-bit Large Language Models (LLMs). In this study, we present \textbf{\text{BitNet b1.58}}, a 1-bit LLM where every parameter is drawn from the ternary set \{-1, 0, 1\}. This architecture delivers performance comparable to full-precision Transformers (e.g., FP16 or BF16) of equivalent model scale and training data—achieving similar perplexity scores and downstream task results. Simultaneously, it delivers substantial reductions in latency, memory footprint, throughput, and energy usage. Beyond these operational advantages, the 1.58-bit LLM introduces a novel scaling law and provides a blueprint for developing future LLMs that excel in both performance and efficiency. Importantly, this work also points toward a new computational paradigm, paving the way for the design of specialized hardware tailored for 1-bit LLM architectures.
\end{abstract}

\section{The Era of 1-bit LLMs}

AI research has recently witnessed unprecedented acceleration in the growth and capability of Large Language Models (LLMs). While these models achieve exceptional outcomes across numerous NLP applications, their ever-increasing size intensifies the complexity of their deployment and exacerbates concerns over high energy demands and related ecological and financial costs.

To mitigate these issues, one promising direction is post-training quantization, which constructs compact, low-bit variants of LLMs specifically for inference~\cite{smoothquant, gptq, quip, quip_sharp}. By lowering the bit-width of model weights and activations, this method markedly cuts down on both memory usage and computational load. Developments have steadily shifted from 16-bit precision toward lower-bit formats, with 4-bit models gaining prominence~\cite{gptq, awq}. Despite its widespread application in industrial settings, post-training quantization often falls short of optimality.

Recent advances in 1-bit model architectures, exemplified by BitNet~\cite{bitnet}, suggest a compelling pathway toward reducing the operational cost of large language models (LLMs) while preserving their effectiveness. Conventional LLMs operate using 16-bit floating-point representations (e.g., FP16 or BF16), with the computational burden dominated by matrix multiplication operations that rely heavily on floating-point addition and multiplication. BitNet, however, transforms this process: its matrix multiplications require only integer addition, dramatically decreasing the energy demands associated with LLM computations. Since the computational throughput of many hardware accelerators is ultimately constrained by available power, these reductions in energy consumption directly enable increases in processing speed.

Beyond computation, inference also incurs substantial expense when moving model parameters from off-chip DRAM to on-chip SRAM accelerators. While enlarging on-chip SRAM is one route to improve throughput, the economic cost is often prohibitive compared to DRAM. 1-bit LLMs address this challenge by offering significantly smaller memory requirements, both in terms of storage capacity and bandwidth usage, compared to their full-precision counterparts. This minimizes both the cost and latency of weight transfers from DRAM, facilitating faster and more efficient inference.

In this paper, we present a noteworthy 1-bit LLM adaptation named \textbf{\text{BitNet b1.58}}, characterized by ternary parameterization with values in the set \{-1, 0, 1\}. By incorporating 0 into the set of possible weights, we extend the original 1-bit BitNet framework to 1.58 bits per parameter under the binary system. \text{BitNet b1.58} retains all computational and memory efficiency gains of 1-bit BitNet, benefiting from a computation paradigm that nearly eliminates multiplication in matrix operations and enables high degrees of optimization. Its power efficiency matches that of the original 1-bit approach, while delivering superior memory efficiency, throughput, and latency compared to FP16 LLM benchmarks. Moreover, \text{BitNet b1.58} introduces two further improvements: firstly, it enhances the model's expressivity through explicit feature filtering enabled by zero-valued weights, which can notably boost 1-bit LLM accuracy; secondly, our experimental results demonstrate that \text{BitNet b1.58} can rival full-precision (FP16) models in perplexity and end-task outcomes, beginning at a model scale of 3B parameters, given identical training regimes and configurations.

\section{BitNet b1.58}

\text{BitNet b1.58} is derived from the BitNet model, which is a Transformer variant that substitutes the standard \emph{nn.Linear} layer with the \emph{BitLinear} layer. This model is trained from the ground up, utilizing weights quantized to 1.58 bits and activations quantized to 8 bits. In relation to the original BitNet, it incorporates several adjustments, which are detailed below.

\paragraph{Quantization Function.} To restrict the possible weight values to the set \{-1, 0, +1\}, we utilize an \emph{absmean} quantization approach. This process begins by normalizing the weight matrix using its mean absolute value and subsequently rounding each entry to its closest integer in \{-1, 0, +1\}:

\begin{equation}
    \widetilde{W} = \mathrm{RoundClip}(\frac{W}{\gamma + \epsilon}, -1, 1),
\end{equation}

\begin{equation}
    \mathrm{RoundClip}(x, a, b) = \max(a, \min(b, \mathrm{round}(x))),
\end{equation}

\begin{equation}
    \gamma = \frac{1}{nm}\sum_{ij} |W_{ij}|.
\end{equation}

For activations, we apply the same quantization method as the original BitNet, with the modification that we avoid rescaling activations before applying the non-linear functions into the interval $[0, Q_b]$. Instead, for each token, activations are normalized to $[-Q_b, Q_b]$, which eliminates the necessity for zero-point quantization. This modification simplifies implementation and benefits system-level optimization, while maintaining similar performance in our empirical evaluations.

\paragraph{LLaMA-alike Components.}

The architecture outlined in LLaMA~\cite{llama, llama2} serves as the default backbone for many community-driven LLMs. To align with open-source standards, the \text{BitNet b1.58} configuration incorporates LLaMA-style modules. Concretely, this involves the use of RMSNorm~\cite{rmsnorm}, SwiGLU~\cite{swiglu}, rotary positional embedding~\cite{rope}, and omitting all bias terms. As a result, \text{BitNet b1.58} is readily compatible with widely-used open-source frameworks (such as Huggingface, vLLM~\cite{vllm}, and llama.cpp\footnote{\url{https://github.com/ggerganov/llama.cpp}}), requiring minimal integration effort.

\section{Results}

We conducted a systematic evaluation of \text{BitNet b1.58} in comparison to a faithfully reproduced FP16 \text{LLaMA LLM} across multiple model sizes. To maintain consistency and fairness, both sets of models underwent pre-training on the RedPajama dataset~\cite{redpajama}, each for a total of 100 billion tokens. Their zero-shot capabilities were assessed across several language benchmarks, including ARC-Easy~\cite{arc}, ARC-Challenge~\cite{arc}, Hellaswag~\cite{hellaswag}, Winogrande~\cite{winoGrande}, PIQA~\cite{piqa}, OpenbookQA~\cite{openbookqa}, and BoolQ~\cite{boolq}. Additionally, we computed validation perplexity metrics on WikiText2~\cite{wikitext2} and C4~\cite{c4} corpora.

In order to analyze inference efficiency, we measured GPU memory consumption and response latency for both \text{LLaMA LLM} and \text{BitNet b1.58}. The evaluations employed the FasterTransformer\footnote{\url{https://github.com/NVIDIA/FasterTransformer}} framework, which is optimized for low-latency LLM inference on GPUs. Notably, the Ladder~\cite{ladder} 2-bit kernel was incorporated for \text{BitNet b1.58}. For all setups, we reported the average time required to generate each output token, as this constitutes the dominant factor in inference runtime.

A summary of the perplexity and efficiency statistics for \text{BitNet b1.58} and \text{LLaMA LLM} is provided in Table~\ref{tab:ppl}. The findings indicate that, beginning with the 3B parameter configuration, \text{BitNet b1.58} achieves comparable perplexity to its full-precision \text{LLaMA LLM} counterpart while offering a 2.71-fold speedup and 3.55-fold reduction in GPU memory footprint. Remarkably, \text{BitNet b1.58} with 3.9B parameters delivers 2.4 times faster inference and utilizes 3.32 times less memory, outperforming the 3B \text{LLaMA LLM} on all tested metrics.

Table~\ref{tab:zero-shot} details the results of zero-shot accuracy on the selected end tasks. Evaluation protocols followed the \emph{lm-evaluation-harness}\footnote{\url{https://github.com/EleutherAI/lm-evaluation-harness}} methodology. These results highlight a diminishing performance gap between \text{BitNet b1.58} and \text{LLaMA LLM} as model scale increases, with \text{BitNet b1.58} matching the full-precision baseline at the 3B parameter scale. Corroborating the perplexity analyses, the end-task scores show \text{BitNet b1.58} 3.9B outperforms \text{LLaMA LLM} 3B, all while requiring less memory and achieving lower latency. Collectively, these observations establish \text{BitNet b1.58} as a Pareto improvement over the leading large language model approaches.

\paragraph{Memory and Latency}

We expanded the model scale to include 7B, 13B, and 70B parameters, and assessed the corresponding costs. As presented in Figure~\ref{fig:latency-memory-energy}, both latency and memory usage exhibit clear scaling patterns, with speed improvements becoming more pronounced as model size increases. Notably, \text{BitNet b1.58} at 70B parameters achieves a 4.1$\times$ acceleration relative to the \text{LLaMA LLM} baseline. This acceleration is attributable to the fact that the runtime for \emph{nn.Linear} operations escalates with growing model dimensions. A parallel pattern is observed in memory requirements: as the embedding layer is retained at full precision, its memory footprint comprises a decreasing fraction of total memory as models become larger. Both latency and memory usage figures were obtained using a 2-bit kernel; hence, further cost reductions are potentially attainable through additional optimization.

\paragraph{Energy}

We further analyze energy consumption by estimating the arithmetic operations energy for both \text{BitNet b1.58} and \text{LLaMA LLM}. The focus is placed on matrix multiplication, the predominant computational cost in LLMs. Figure~\ref{fig:energy} depicts the breakdown of energy expenditure, revealing that \text{BitNet b1.58} primarily involves INT8 additions, whereas \text{LLaMA LLM} utilizes both FP16 additions and multiplications. Based on the energy modeling from~\cite{energycost, pokebnn}, \text{BitNet b1.58} demonstrates a 71.4-fold reduction in arithmetic operations energy for matrix multiplications when deployed on 7nm chips. Furthermore, we present the end-to-end energy usage for sequences consisting of 512 tokens. The results indicate that increasing the model size yields enhanced energy efficiency for \text{BitNet b1.58} compared to the FP16 \text{LLaMA LLM} baseline. This can be attributed to the rising proportion of computation spent in \emph{nn.Linear} operations for larger models, while the relative overhead from ancillary components diminishes as model size grows.

\paragraph{Throughput}

We assess the throughput capabilities of \text{BitNet b1.58} and \text{LLaMA LLM} with 70B parameters, deployed on a pair of 80GB A100 GPUs and utilizing pipeline parallelism~\cite{gpipe} to enable execution of the \text{LLaMA LLM} 70B model. The batch size was incrementally increased until the GPUs' memory was fully utilized, using a fixed sequence length of 512. According to Table~\ref{tab:throughput}, \text{BitNet b1.58} 70B accommodates up to 11-fold greater batch sizes relative to \text{LLaMA LLM}, thereby delivering throughput that is 8.9 times higher.

\textbf{\text{BitNet b1.58} establishes a novel scaling law relating model quality to inference efficiency.} To contextualize this, equivalences between various 1.58-bit and 16-bit model sizes can be established based on the experimental results presented in Figure~\ref{fig:latency-memory-energy} and \ref{fig:energy}.

\begin{itemize}
    \item The 13B BitNet b1.58 outperforms the 3B FP16 LLM with respect to latency, memory footprint, and energy requirements.
    \item The 30B BitNet b1.58 demonstrates greater efficiency in latency, memory consumption, and energy usage compared to the 7B FP16 LLM.
    \item The 70B BitNet b1.58 is more resource-efficient—regarding latency, memory, and energy—than the 13B FP16 LLM.
\end{itemize}

\paragraph{Training with 2T Tokens}

The total number of tokens used in training is a pivotal consideration for LLM performance. To examine how well \text{BitNet b1.58} scales with token count, we trained a \text{BitNet b1.58} model using 2T tokens, adhering to the data preparation approach of StableLM-3B~\cite{StableLM-3B-4E1T}, the leading open-source 3B parameter model. We evaluated both models on a composite benchmark, incorporating Winogrande~\cite{winoGrande}, PIQA~\cite{piqa}, SciQ~\cite{sciq}, LAMBADA~\cite{lambada}, and ARC-easy~\cite{arc}. Table~\ref{tab:2t} displays the reported zero-shot accuracies. For evaluation metrics that include both accuracy and normalized accuracy, we computed the mean. The 2T token results for StableLM 3b are sourced from its official technical documentation. Our analysis reveals that \text{BitNet b1.58} consistently outperforms across all tasks, demonstrating that 1.58-bit LLMs maintain robust generalization performance.

\section{Discussion and Future Work}

\textbf{1-bit Mixture-of-Experts (MoE) LLMs}

Mixture-of-Experts (MoE) architectures have emerged as an efficient solution for scaling LLMs, notably decreasing the computational requirements measured in FLOPs. Nonetheless, MoE models often face constraints due to their substantial memory demands and the significant inter-device communication necessary for their deployment. The advent of 1.58-bit LLMs offers a way to overcome these limitations. The minimized memory usage not only decreases the hardware needed for MoE model deployment, but also substantially lessens the communication bandwidth necessary for sharing activations among devices. Ideally, this reduction can enable the entire model to reside within the confines of a single chip, effectively eliminating cross-chip communication overhead altogether.

\textbf{Native Support of Long Sequence in LLMs}

Supporting long input sequences is increasingly essential for LLM applications. A persistent obstacle for such long sequence inference is the memory cost arising from KV cache storage. \text{BitNet b1.58} marks a substantial improvement in addressing this concern, shrinking the bit-width of activations from 16 to 8 and thereby doubling the attainable context length using the same memory capacity. There exists further potential to compress activations losslessly down to 4 bits or less within the 1.58-bit LLM framework, which we propose to investigate in future work.

\textbf{LLMs on Edge and Mobile}

1.58-bit LLMs present compelling advantages for deployment on edge and mobile platforms, where resources such as memory and computational throughput are inherently restricted. By reducing both the energy and memory requirements, 1.58-bit LLMs can be more feasibly implemented on such constrained devices, thereby unlocking new classes of language-based applications that were previously infeasible. These advances not only broaden the functionality available on mobile and edge hardware but also take advantage of CPUs—the predominant processing units in these contexts—due to the compatibility of 1.58-bit LLMs with CPU architectures. Consequently, \text{BitNet b1.58} is particularly well-suited for efficient execution on these devices, driving further improvements in performance and capabilities.

\textbf{New Hardware for 1-bit LLMs}

There has been notable progress, such as in the Groq\footnote{\url{https://groq.com/}} project, toward developing custom hardware (e.g., LPUs) dedicated to supporting LLMs. In light of the computational paradigm shift ushered in by BitNet~\cite{bitnet}, we advocate for the creation of new hardware and systems optimized for 1-bit LLMs, encouraging the research community to innovate in this direction.